% Authoritative LaTeX chapter. Edit this file for future revisions.
% Initially migrated 2026-09-11; no Markdown import occurs during PDF builds.
\chapter{Installation and First Run}
\label{chap:02_installation}
Contents · \hyperref[chap:03_options]{Next: C options}

\section{Components and prerequisites}
\label{sec:auto:02_installation:1}
The recommended starting point is the C reference implementation. It requires a C11-capable GCC, OpenMP, and Git. The repository's Linux guidance specifies GCC 11 or later; Windows uses MinGW-w64 GCC. See \hyperref[chap:06_python_batch]{Appendix~\ref*{chap:06_python_batch}} for the Python batch package's dependencies and supported scope.


{\TableFont
\begin{longtable}{@{}>{\raggedright\arraybackslash}p{\dimexpr 0.3400\linewidth-4.00pt\relax}>{\raggedright\arraybackslash}p{\dimexpr 0.6600\linewidth-4.00pt\relax}@{}}
\toprule
\textbf{Location} & \textbf{Role} \\
\midrule
\endfirsthead
\toprule
\textbf{Location} & \textbf{Role} \\
\midrule
\endhead
\bottomrule
\endfoot
\texttt{code/\allowbreak{}OCRT\_\allowbreak{}C/\allowbreak{}src/\allowbreak{}} & C source \\
\texttt{code/\allowbreak{}OCRT\_\allowbreak{}C/\allowbreak{}inputs/\allowbreak{}} & AFGL, absorption cross sections, seawater IOPs, and Mie data read by C \\
\texttt{code/\allowbreak{}OCRT\_\allowbreak{}C/\allowbreak{}build/\allowbreak{}} & Executables built locally \\
\texttt{code/\allowbreak{}OCRT\_\allowbreak{}Python/\allowbreak{}} & Python implementation \\
\texttt{scripts/\allowbreak{}} & Public Mie data installers and hash lists \\
\end{longtable}
}

The command examples assume that the repository is cloned into a folder named \texttt{OCRT}. They do not hard-code absolute paths that vary by account or computer.


\par\noindent\begin{minipage}{\linewidth}
\begin{ManualCode}
git clone https://github.com/brtnt/OCRT.git
cd OCRT
git rev-parse HEAD
\end{ManualCode}
\end{minipage}\par

If the repository is already cloned, inspect any changes to the existing source and decide which version to use. Executables are not distributed through Git, so build after obtaining the source.

\section{Installing input data}
\label{sec:auto:02_installation:2}
Small data files are included in Git. The 181 FR631 Mie files needed for aerosol and particulate-ocean calculations are distributed in four archives in the \href{https://github.com/brtnt/OCRT/releases/tag/data-v1}{data-v1 release}. Allow approximately 1.3 GB for the download, approximately 4.3 GB for the C inputs, and approximately 4.3 GB for the Python copy. Also reserve space for the archive cache.

Run from the repository root.

Linux/Bash:


\par\noindent\begin{minipage}{\linewidth}
\begin{ManualCode}
bash scripts/fetch_data.sh
\end{ManualCode}
\end{minipage}\par

Windows/PowerShell:


\par\noindent\begin{minipage}{\linewidth}
\begin{ManualCode}
powershell -ExecutionPolicy Bypass -File scripts\fetch_data.ps1
\end{ManualCode}
\end{minipage}\par

The scripts verify the archives' SHA-256 hashes and install the data in the C and Python input folders. Do not use \texttt{-\allowbreak{}NoVerify} for a normal installation. Existing input files with the same names may be updated during installation, so keep modified user data separately.

To verify individual C Mie files on Linux as well:


\par\noindent\begin{minipage}{\linewidth}
\begin{ManualCode}
cd code/OCRT_C/inputs
sha256sum -c ../../../scripts/MIE_SHA256SUMS_data-v1.txt
cd ../../..
\end{ManualCode}
\end{minipage}\par

C Rayleigh examples without aerosols or underwater-particle calculations can be run without downloading the large Mie data set. Calculations with gas absorption enabled require the small AFGL and cross-section files. \textbf{Python has paths that preload additional data in constructors, so the C conditions for running without those data do not transfer directly to Python.}

\section{Building on Linux}
\label{sec:auto:02_installation:3}
The default CPU option in the \texttt{Makefile} is \texttt{-\allowbreak{}march=\allowbreak{}cascadelake} (AVX-512). Execution fails if the CPU does not support it. For a typical x86-64 AVX2 environment, specify the options explicitly:


\par\noindent\begin{minipage}{\linewidth}
\begin{ManualCode}
cd code/OCRT_C
make CFLAGS='-std=c11 -O3 -march=x86-64-v3 -ffp-contract=fast -fassociative-math -fno-signed-zeros -fno-trapping-math -DOCRT_FAST_KERNELS -fopenmp -Isrc'
./build/ocrt --version
./build/ocrt --help
\end{ManualCode}
\end{minipage}\par

When rebuilding with a different CPU option, force a rebuild with the same \textbf{complete} CFLAGS, for example \texttt{make -\allowbreak{}B CFLAGS=\allowbreak{}'\ldots{}'}. Replace \texttt{\ldots{}} in that command with the actual options. Make does not automatically regenerate an existing executable merely because CFLAGS changed.

An x86-64 CPU without AVX2 can use a build with \texttt{-\allowbreak{}march=\allowbreak{}x86-\allowbreak{}64}, but performance and numerical differences must be checked again. The x86-specific options above cannot be used on ARM/macOS, which is outside this manual's verified execution scope. Do not treat Apple Clang as equivalent to the GCC/OpenMP environment.

\section{Building on Windows}
\label{sec:auto:02_installation:4}
Install a GCC environment following the \href{https://www.msys2.org/}{official MSYS2 installation guidance}. The current official default environment is UCRT64. In a UCRT64 terminal, install GCC with:


\par\noindent\begin{minipage}{\linewidth}
\begin{ManualCode}
pacman -S mingw-w64-ucrt-x86_64-gcc
gcc --version
\end{ManualCode}
\end{minipage}\par

Existing OCRT Windows scripts were written for MinGW64 GCC. Consistently use the \texttt{bin} path of your MinGW64 or UCRT64 environment, and do not mix GCC/DLL environments. The following is \textbf{an example for UCRT64 with MSYS2 installed in its default location}. If your installation is elsewhere, change the corresponding PATH entry.

In PowerShell, starting at the repository root:


\par\noindent\begin{minipage}{\linewidth}
\begin{ManualCode}
$env:Path = 'C:\msys64\ucrt64\bin;' + $env:Path
gcc --version
Set-Location code\OCRT_C
.\build_win.bat
.\build\ocrt_x86-64-v3.exe --version
.\build\ocrt_x86-64-v3.exe --help
\end{ManualCode}
\end{minipage}\par

\texttt{build\_\allowbreak{}win.\allowbreak{}bat} includes the Windows compatibility source files, and its default output is \texttt{build\textbackslash{}ocrt\_\allowbreak{}x86-\allowbreak{}64-\allowbreak{}v3.\allowbreak{}exe}. It waits for a keypress after running. \texttt{build\_\allowbreak{}win.\allowbreak{}bat native} targets the current computer's CPU. Use \texttt{cascadelake} only after confirming AVX-512 support. An old \texttt{ocrt.\allowbreak{}exe} remaining in the folder is not automatically replaced by the new build.

The \texttt{build/\allowbreak{}run\_\allowbreak{}star.\allowbreak{}py} mentioned in the script's final message is not a required distribution file in the public repository. Its absence does not indicate a compilation failure. Verify the build using the first calculation in Appendix~\ref{sec:example_first_run} and the validation procedure in Appendix~\ref{sec:example_ipss_gates}.

This manual verified building and running on Linux/WSL. Native Windows and UCRT64 builds require a separate environment check.

\section{Working directory at runtime}
\label{sec:auto:02_installation:5}
\textbf{Set the current working directory to \texttt{code/\allowbreak{}OCRT\_\allowbreak{}C} when running.} Specifying the executable's location alone does not automatically change the internal \texttt{inputs/\allowbreak{}} relative paths.

Bash and PowerShell commands for a single calculation, with output guidance, are provided in Appendix~\ref{sec:example_first_run}. For a first installation, begin with the Quick start at the front of the manual.

\section{Validation and reproducibility}
\label{sec:auto:02_installation:6}
Geometry and Rayleigh numerical checks can be run after installation. The validation script and its completion markers are described in Appendix~\ref{sec:example_ipss_gates}.

With the outputs, retain the commit, build command and GCC version, CPU, executable hash, input-data version, complete options, environment variables, output headers, and standard-error logs. Perform scientific comparisons under matching conditions, distinguishing last-digit differences caused by the compiler/CPU from model differences.

Also pay attention to case sensitivity, paths containing spaces, and UTF-8 handling in inputs and outputs. Save input CSV files as UTF-8 without a BOM, with \texttt{.\allowbreak{}} as the decimal separator and \texttt{,\allowbreak{}} as the field delimiter.

